\section{The first retained class on the regular set}
\label{p14:comp:stopped}

This section specializes the general stopping identity of
Section~\ref{p1:stop:section} to the two leading pair channels. It
supplies the boundary estimates used by the first reconstruction and
the positional comparison. The optimal scalar reconstruction uses the
larger class $\mathscr D_2$ in Section~\ref{p13:transfer:section}.
Here and in the regular-set reconstruction below, write
\[
 G=G_N,\quad n=|G|,\quad Z_G=Z_N^\circ,\quad
 \mu=np,\quad E_N=\mathcal E_N,
 \qquad a_N=\E(Z_N)_2-\lambda_N^2.
\]
The pair calculation and mean estimate give
$a_N-\alpha_N=O_{C,\varepsilon}(N^{-1+\varepsilon})$ and
$a_N\asymp_C LN^{-2/3}>0$ for large $N$.
Every space on $G$ has dimension $L$, and every marginal is $p$.

\subsection{Retaining a direct family and one compressed pair}

We now use the regular set $G$ of arithmetic windows.
For $I\subseteq G$, recall
\[
 \rho(I)=|I|L-\dim\sum_{x\in I}S_x.
\]
Every $S_x$ has dimension $L$, and a family is called direct
when $\rho(I)=0$.  Let $E_N$ be the family of good pairs
from Corollary~\ref{p1:arith:edgefamily}.  Its pairs are
compatible and dependent, and
\begin{equation}
 \rho(e)\le L/3+O(1)\qquad(e\in E_N).
 \label{p1:stop:pair-nullity}
\end{equation}
For $e=\{x,y\}$ put $V_e=S_x+S_y$ and $q_e=q(e)$.  Define
\begin{equation}
 \Gamma_e(h)=\left\{
 U\subseteq G\setminus e:|U|=h,\quad
 V_e,\ (S_z)_{z\in U}\ \text{are in direct sum}
                    \right\}.
 \label{p1:stop:gamma}
\end{equation}
The direct-sum condition concerns the single compressed block
$V_e$ and the other blocks.  It does not assert independence
of the two spaces inside $e$.  The affine probability formula
gives
\begin{equation}
 q(e\cup U)=q_ep^{|U|}
              \qquad\bigl(U\in\Gamma_e(|U|)\bigr).
 \label{p1:stop:direct-pair-probability}
\end{equation}

Let $\mathscr D$ consist of all direct families together with
all families $e\cup U$ occurring in
\eqref{p1:stop:gamma}.  Thus a family in $\mathscr D$ is allowed
one good pair, treated as a single block, but no further dependence.

\begin{lemma}[Heredity and boundary nullity]
\label{p1:stop:class}
The family $\mathscr D$ is hereditary.  Each non-direct member
contains a unique dependent pair $e\in E_N$, and its nullity
equals $\rho(e)$.  Hence, for all sufficiently large $N$,
\begin{equation}
 \partial\mathscr D
 \subseteq\{I\notin\mathscr D:
                 0<\rho(I)\le 4L/3+O(1)<3L/2\}.
 \label{p1:stop:boundary-nullity}
\end{equation}
\end{lemma}

\begin{proof}
Every subfamily of a direct family is direct.
Consider $e\cup U$ in the second part of the definition.
A subfamily which retains both endpoints of $e$ still consists
of $V_e$ and some of the independent singleton blocks.  A subfamily
which removes an endpoint is direct, since the remaining endpoint
space is contained in $V_e$.  This proves heredity.

The direct-sum condition excludes a dependent pair inside $U$
and excludes a dependent pair joining $U$ to either endpoint of
$e$.  Thus $e$ is the unique dependent pair.  Additivity of the
ranks of the compressed blocks then gives
$\rho(e\cup U)=\rho(e)$.

Adding one $L$-dimensional space increases nullity by at most
$L$.  The proper prefix of a boundary member has nullity at most
$L/3+O(1)$, so the member itself has nullity at most
$4L/3+O(1)$.  Finally, nullity zero would make it a direct
family and therefore a member of $\mathscr D$.
\end{proof}

The retained dependence has the size of the principal correction.
The strict upper bound on the boundary nullity is what makes
the stopped remainder easier to estimate than the unrestricted
absolute expansion.

\subsection{The factorial tail of the boundary}

Write $n=|G|$ and $\mu=np$, and choose a fixed upper bound
$\Lambda$ for $\mu$ throughout the critical band.  We need a
cutoff large enough to absorb the loss $2^{3L/2}$ in
\eqref{p1:stop:boundary-nullity}, yet small enough for the
critical-size component estimate.  The following explicit
choice meets both requirements:
\begin{equation}
 K_N=\min\left\{K\in\mathbb N:
 K\ge\max(4,\lceil4e\Lambda\rceil),\quad
 K\log\frac{K}{2e\Lambda}\ge3\log N\right\}.
 \label{p1:stop:cutoff}
\end{equation}
It satisfies
\begin{equation}
 K_N=(3+o(1))\frac{\log N}{\log\log N}=o(L);
 \qquad K_N\le L/4
 \quad\text{for all sufficiently large }N.
 \label{p1:stop:cutoff-range}
\end{equation}
The coefficient $3$ is a convenient sufficient choice.

\begin{lemma}[Large boundary families]
\label{p1:stop:tail}
Uniformly in the fixed critical band,
\begin{equation}
 \sum_{\substack{I\in\partial\mathscr D\\|I|>K_N}}
           2^{|I|}q(I)=O_C(N^{-3/2}).
 \label{p1:stop:tail-bound}
\end{equation}
\end{lemma}

\begin{proof}
The affine law and \eqref{p1:stop:boundary-nullity} imply
\[
 q(I)\le 2^{3L/2}p^{|I|}
          \ll_C N^{3/2}p^{|I|}.
\]
This remains valid for incompatible families, whose probability
is zero.  Therefore
\[
 \sum_{\substack{I\in\partial\mathscr D\\|I|>K_N}}
       2^{|I|}q(I)
 \ll_C N^{3/2}
       \sum_{r>K_N}\binom nr(2p)^r
 \le C_C N^{3/2}
       \sum_{r>K_N}\frac{(2\mu)^r}{r!}.
\]
By $K!\ge(K/e)^K$ and \eqref{p1:stop:cutoff},
$(2\mu)^{K_N}/K_N!\le N^{-3}$.
The ratios of subsequent terms are at most
$2\Lambda/(K_N+1)<1/2$, so the last tail is at most
$2N^{-3}$.  This proves the assertion.
\end{proof}

Only families of size at most $K_N$ remain.  They have nullity
at most $3L/2$, a bound fixed relative to $L$.  The component
estimate in Theorem~\ref{p1:core:bounded-nullity} supplies the
arithmetic input for this range.  Section~\ref{p1:close:section}
will convert that centered estimate into the positive boundary
mass required by \eqref{p1:stop:norm}.
